\documentclass{article}

\usepackage{amsmath, amsfonts, amssymb, amsthm}
\usepackage{parskip}

\newtheorem*{claim}{Claim}

\title{Proof by Induction: An Integral with a Parameter}
\author{Introduction to Engineering Mathematics}
\date{November 2026}

\begin{document}
\maketitle

\begin{claim}
  For every integer $n \ge 0$,
  \[
    I_n = \int_0^1 (1 - x^2)^n \, dx = \frac{2^{2n} (n!)^2}{(2n + 1)!}.
  \]
\end{claim}

\section*{Preliminary: a reduction formula}

For $n \ge 1$, integrate by parts with
\[
  u = (1 - x^2)^n, \quad dv = dx,
  \qquad
  du = -2nx (1 - x^2)^{n-1} \, dx, \quad v = x.
\]
This gives
\[
  I_n = \Big[ x (1 - x^2)^n \Big]_0^1 + 2n \int_0^1 x^2 (1 - x^2)^{n-1} \, dx.
\]
The boundary term is zero, since $x (1 - x^2)^n$ vanishes at both $x = 0$ and
$x = 1$. In the remaining integral, write $x^2 = 1 - (1 - x^2)$:
\[
  I_n = 2n \int_0^1 \left[ (1 - x^2)^{n-1} - (1 - x^2)^n \right] dx
      = 2n \left( I_{n-1} - I_n \right).
\]
Solving for $I_n$ gives the reduction formula
\[
  I_n = \frac{2n}{2n + 1} \, I_{n-1}, \qquad n \ge 1.
\]

\section*{Proof outline}

Let $P(n)$ denote the statement ``$I_n = \dfrac{2^{2n} (n!)^2}{(2n + 1)!}$''.
We show that $P(n)$ holds for all $n \ge 0$ by induction on $n$.

\subsection*{Step 1: Base case ($n = 0$)}

For $n = 0$ the integrand is $1$, so
\[
  I_0 = \int_0^1 1 \, dx = 1.
\]
On the other hand,
\[
  \frac{2^0 (0!)^2}{1!} = 1.
\]
Both sides agree, so $P(0)$ holds.

\subsection*{Step 2: Induction hypothesis}

Assume that $P(k)$ holds for some integer $k \ge 0$, that is,
\[
  I_k = \frac{2^{2k} (k!)^2}{(2k + 1)!}.
\]

\subsection*{Step 3: Induction step}

We must show that $P(k + 1)$ holds, i.e.\ that
\[
  I_{k+1} = \frac{2^{2k+2} ((k + 1)!)^2}{(2k + 3)!}.
\]
By the reduction formula with $n = k + 1$, and then the induction hypothesis,
\[
  I_{k+1} = \frac{2k + 2}{2k + 3} \, I_k
          = \frac{2k + 2}{2k + 3} \cdot \frac{2^{2k} (k!)^2}{(2k + 1)!}.
\]
To bring this into the required form, multiply numerator and denominator by
$2k + 2$:
\[
  I_{k+1} = \frac{(2k + 2)^2 \, 2^{2k} (k!)^2}{(2k + 3)(2k + 2)(2k + 1)!}
          = \frac{2^2 (k + 1)^2 \, 2^{2k} (k!)^2}{(2k + 3)!}
          = \frac{2^{2k+2} ((k + 1)!)^2}{(2k + 3)!},
\]
using $(2k + 3)(2k + 2)(2k + 1)! = (2k + 3)!$ and $(k + 1) \, k! = (k + 1)!$.
This is exactly $P(k + 1)$.

\subsection*{Step 4: Conclusion}

Since $P(0)$ is true, and $P(k)$ implies $P(k + 1)$ for every $k \ge 0$, the
principle of mathematical induction tells us that $P(n)$ is true for all
$n \ge 0$. \qed

\end{document}
